\chapter{Lezione 1b - }

\section{Fisica dei Plasmi e Teoria della Dinamo: Regimi e Ionizzazione}

Il materiale che segue costituisce gli appunti relativi al Modulo 1b del corso di Fisica dei Plasmi e Teoria della Dinamo (tenuto dal Prof. Alfio Bonanno presso l'INAF – Osservatorio Astrofisico di Catania). L'oggetto principale di questa trattazione verte sullo studio dei regimi fisici dello stato di plasma e sull'analisi termodinamica dei processi di ionizzazione.

\subsection{Programma e Riferimenti Bibliografici}
Il contenuto della presente lezione è strutturato secondo due direttrici principali:
\begin{enumerate}
\item L'analisi del parametro di plasma e la classificazione dei vari regimi fisici in cui un plasma può trovarsi.
\item Lo studio dell'equilibrio di ionizzazione, formulato matematicamente attraverso l'Equazione di Saha.
\end{enumerate}

Per completare la preparazione su questi argomenti, si raccomanda la lettura dei seguenti testi di riferimento: Boyd \& Sanderson (capitolo 1), Bittencourt (capitolo 1), Fitzpatrick (capitolo 1) e Zel'dovich \& Raizer (capitolo III). Lo studio va integrato con le note del corso (capitolo 1, paragrafi 1.4-1.5) e con lo svolgimento pratico del Set 1 di problemi.

\subsection{Riepilogo dei Concetti Precedenti}
Prima di addentrarci nei nuovi concetti, è essenziale richiamare in modo sistematico i risultati teorici chiave derivati nel modulo precedente, i quali fungono da base concettuale per le argomentazioni odierne. Tali capisaldi includono:
\begin{itemize}
\item La \textbf{lunghezza di Debye} $\lambda_D = (T / 4\pi n e^2)^{1/2}$, che quantifica la scala spaziale su cui si manifesta lo schermo elettrostatico e garantisce l'assunzione di quasi-neutralità macroscopica del mezzo.
\item La \textbf{frequenza di plasma} (elettronica) $\omega_{pe} = (4\pi n e^2 / m_e)^{1/2}$, che descrive la pulsazione naturale delle oscillazioni collettive delle cariche e si lega alla lunghezza di Debye tramite la velocità termica: $\omega_{pe}\lambda_D = v_{te}$.
\item L'\textbf{energia di correlazione} termodinamica $\vert{}w\vert{}/T = 1 / 12N_D$. Questa relazione dimostra che il plasma obbedisce all'approssimazione di gas ideale esclusivamente quando il numero di particelle racchiuse in una sfera di Debye risulta essere molto elevato ($N_D \gg 1$).
\end{itemize}
Forte di queste premesse, l'obiettivo odierno è determinare con precisione sotto quali condizioni e in quali domini parametrici un plasma possa essere considerato \textbf{ideale}, governato dalla meccanica \textbf{classica} (non degenere) e adeguatamente \textbf{ionizzato}.

\section{Parametro di Plasma e Regimi Fisici}

In questa sezione esploriamo i criteri per definire e classificare i vari regimi fisici in cui un plasma può esistere, analizzando il grado di accoppiamento e i limiti della meccanica classica rispetto alla quantistica.

\subsection{Energia Potenziale vs Energia Cinetica: Il Parametro di Accoppiamento}
Per determinare se il comportamento di un plasma sia dominato dalle interazioni a due corpi o dall'agitazione termica collettiva, si introduce il concetto di \textbf{raggio di Wigner-Seitz} ($a$). Questo raggio rappresenta la distanza media tra due particelle adiacenti nel mezzo, calcolata assumendo che ogni particella occupi un volume sferico esclusivo. Si ricava dall'identità $\frac{4\pi}{3} a^3 = \frac{1}{n}$, da cui:
\begin{equation}
a = \left( \frac{3}{4\pi n} \right)^{1/3}
\end{equation}
dove $n$ è la densità numerica delle particelle.

Utilizzando questa scala di distanza spaziale, possiamo definire il \textbf{parametro di accoppiamento} ($\Gamma$). Questo parametro adimensionale è costruito come il rapporto tra l'energia elettrostatica (Coulombiana) di interazione tra particelle vicine e l'energia termica caratteristica del mezzo ($T$):
\begin{equation}
\Gamma \equiv \frac{\text{Energia di Coulomb tra particelle vicine}}{\text{Energia termica}} = \frac{Z^2 e^2}{a T}
\end{equation}

Per evidenziare la connessione tra questo nuovo parametro e le quantità derivate in precedenza, consideriamo un plasma ideale a singola carica ionica ($Z=1$). Sostituiamo le definizioni del quadrato della lunghezza di Debye ($\lambda_D^2 = T / 4\pi n e^2$) e del cubo del raggio di Wigner-Seitz ($a^3 = 3 / 4\pi n$).
Elevando al quadrato il parametro di plasma $N_D$, si ottiene:
\begin{equation}
N_D^2 = \left( \frac{4\pi}{3} \right)^2 n^2 \lambda_D^6 = \frac{T^3}{36\pi n e^6}
\end{equation}
Parallelamente, elevando al cubo il parametro di accoppiamento $\Gamma$, risulta:
\begin{equation}
\Gamma^3 = \frac{e^6}{a^3 T^3} = \frac{4\pi n e^6}{3 T^3}
\end{equation}
Moltiplicando queste due espressioni, tutti i termini fisici (densità, temperatura, carica elementare) si elidono perfettamente, restituendo un prodotto puramente numerico:
\begin{equation}
N_D^2 \Gamma^3 = \frac{1}{27}
\end{equation}
Isolando $N_D$, giungiamo al \textbf{risultato chiave} che lega indissolubilmente il parametro di plasma alla forza di accoppiamento termodinamico:
\begin{equation}
N_D = \frac{1}{3\sqrt{3} , \Gamma^{3/2}}
\end{equation}

\subsection{Regimi di Accoppiamento Debole e Forte}
L'espressione appena derivata rivela l'esistenza di due regimi fisici profondamente distinti, determinati dal valore assunto da $\Gamma$ (o inversamente da $N_D$).

\textbf{1. Plasma debolmente accoppiato (Weakly coupled):}
Si verifica quando l'energia termica sovrasta le interazioni di vicinato, ovvero $\Gamma \ll 1$. Per la formula precedentemente derivata, ciò implica necessariamente $N_D \gg 1$. In questo regime:
\begin{itemize}
\item Ogni singola carica è schermata in modo efficiente dall'azione collettiva di \textit{molte} particelle circostanti racchiuse nella nube di Debye.
\item L'evoluzione dinamica è guidata da campi elettromagnetici "lisci" (smooth) e macroscopicamente mediati, descrivibili dalle equazioni di Vlasov o della fluidodinamica, permettendo al sistema di obbedire all'equazione di stato (EOS) del \textbf{gas ideale}.
\item I trasferimenti di impulso cinetico sono dominati da innumerevoli collisioni Coulombiane a \textit{piccolo angolo}, mentre gli urti violenti a due corpi (large-angle collisions) sono statisticamente rari.
\end{itemize}
Considerando che $N_D \propto T^{3/2} / n^{1/2}$, si deduce che l'idealità del plasma è massimizzata in condizioni in cui il gas si presenta simultaneamente \textbf{estremamente caldo e molto diluito}. \textit{(Nota convenzionale: nella letteratura scientifica si incontrano diverse definizioni per quantificare la molteplicità particellare; alcuni autori usano $n\lambda_D^3$, altri il parametro $N_D$, e altri ancora pongono l'argomento del Logaritmo di Coulomb pari a $\Lambda = 12\pi n \lambda_D^3$)}.

\textbf{2. Plasma fortemente accoppiato (Strongly coupled):}
Si verifica nel limite opposto, quando le forze inter-particellari competono con l'agitazione termica ($\Gamma \gtrsim 1$), abbattendo il numero di cariche per sfera di Debye ($N_D \sim 1$ o inferiore). In queste condizioni di gas freddo e iperdenso:
\begin{itemize}
\item Le fluttuazioni di carica a corto raggio dominano, introducendo \textbf{correlazioni strutturali} paragonabili a quelle riscontrabili nei liquidi.
\item L'intero impianto teorico dello \textbf{schermo di Debye fallisce}, in quanto decade l'assunzione di un campo di sfondo continuo e linearizzabile.
\item Se si spinge il parametro oltre la soglia critica di $\Gamma \approx 175$ (valore di riferimento per un plasma a singola componente), l'energia cinetica collassa del tutto e il plasma subisce un processo di transizione di fase, arrivando alla \textbf{cristallizzazione} della struttura (reticolo di Wigner).
\end{itemize}

\subsection{Degenerazione Quantistica: I Limiti della Meccanica Classica}
Oltre al grado di accoppiamento elettrostatico, è vitale stabilire a partire da quali soglie fisiche sia imperativo abbandonare l'approccio classico (gas di Maxwell-Boltzmann) per adottare la meccanica quantistica (statistica di Fermi-Dirac).
Gli effetti quantistici non possono più essere trascurati quando l'estensione probabilistica associata alla particella — rappresentata dalla \textbf{lunghezza d'onda termica di de Broglie} ($\lambda_{\text{th}}$) — si dilata fino a diventare confrontabile con la spaziatura media tra le particelle stesse. Essendo $\lambda_{\text{th}} \propto m^{-1/2}$, gli elettroni (avendo una massa a riposo quasi 2000 volte inferiore ai protoni) manifestano la degenerazione molto prima degli ioni.
La condizione per l'insorgenza degli effetti quantistici elettronici si esprime matematicamente con la disuguaglianza:
\begin{equation}
n \lambda_{\text{th}}^3 \gtrsim 1, \quad \text{con} \quad \lambda_{\text{th}} = \frac{h}{(2\pi m_e T)^{1/2}}
\end{equation}
.

Per quantificare il livello di saturazione degli stati quantici, analizziamo il limite in cui la temperatura si approssima allo zero assoluto ($T=0$). Per il Principio di Esclusione di Pauli, i fermioni (elettroni) riempiono progressivamente i livelli energetici dello spazio degli impulsi dal minimo fino a un valore massimo, definito come \textbf{impulso di Fermi} ($p_F$). Considerando che ogni cella elementare dello spazio delle fasi (di volume $(2\pi\hbar)^3$) può ospitare due elettroni con spin opposto, la densità spaziale è legata al volume della "sfera di Fermi" nello spazio degli impulsi dalla relazione:
\begin{equation}
n_e = \frac{2}{(2\pi\hbar)^3} \cdot \frac{4\pi}{3} p_F^3
\end{equation}
. Invertendo l'equazione, si ricava l'espressione per l'impulso di Fermi in funzione della densità di elettroni:
\begin{equation}
p_F = \hbar(3\pi^2 n_e)^{1/3}
\end{equation}
.

Da questo impulso si deduce l'energia cinetica massima posseduta dalle particelle a temperatura nulla, che permette di definire il \textbf{risultato chiave} di questa trattazione, ovvero la \textbf{Temperatura di Fermi} ($T_F$). Esprimendola in unità di energia, essa vale:
\begin{equation}
T_F = \frac{p_F^2}{2m_e} = \frac{\hbar^2 (3\pi^2 n_e)^{2/3}}{2m_e} = 4.23 \times 10^{-11} \left( \frac{n_e}{\text{cm}^{-3}} \right)^{2/3} \text{ K}
\end{equation}
.
Questo parametro agisce da spartiacque termodinamico: un plasma si definisce \textbf{classico} fintanto che la sua temperatura termica $T$ è significativamente maggiore della temperatura di Fermi del gas elettronico ($T \gg T_F$). Qualora si verifichi il caso opposto ($T \lesssim T_F$), la pressione del sistema non è più sorretta dall'agitazione termica, bensì dal principio di Pauli, e il plasma diventa \textbf{degenere}.

\section{Accoppiamento nei Plasmi Degeneri}

Quando la densità del plasma aumenta a tal punto da innescare la degenerazione quantistica, l'analisi termodinamica subisce un profondo mutamento concettuale. Come già stabilito, un plasma è ben descritto dalla fisica classica se la sua temperatura termica $T$ è nettamente superiore alla temperatura di Fermi $T_F$ ($T \gg T_F$). Tuttavia, quando il sistema si raffredda o si comprime fino a raggiungere il regime $T \ll T_F$, l'energia cinetica caratteristica dominante non è più dettata dalla temperatura $T$, bensì dall'energia di Fermi $E_F$.

Per valutare il parametro di accoppiamento in questo regime quantistico, dobbiamo rapportare l'energia potenziale Coulombiana all'energia di Fermi. Sapendo che $n^{2/3} \sim a^{-2}$ (essendo $a$ il raggio interparticellare) e che l'energia di Fermi scala come $E_F \sim \hbar^2 / m_e a^2$, possiamo impostare la seguente stima adimensionale per il parametro di accoppiamento degenere (indicato con $r_s$):
\begin{equation}
\frac{e^2/a}{E_F} \sim \frac{e^2/a}{\hbar^2 / m_e a^2} = \frac{m_e e^2 a}{\hbar^2} = \frac{a}{a_B} \equiv r_s
\end{equation}
In questa formulazione, $a_B$ rappresenta il \textbf{raggio di Bohr} fondamentale dell'atomo di idrogeno, pari a $a_B = \hbar^2 / m_e e^2 \approx 0.53 \text{ \AA}$.

Questa derivazione conduce a un \textbf{risultato chiave} controintuitivo: al contrario dei gas classici, \textbf{i plasmi degeneri diventano sempre più ideali man mano che vengono compressi}. La condizione di idealità quantistica ($r_s < 1$) si verifica per densità elettroniche formidabili, calcolabili imponendo la disuguaglianza $n_e > 3/4\pi a_B^3$, che restituisce la soglia di circa $1.6 \times 10^{24} \text{ cm}^{-3}$.
Per riassumere e confrontare le dinamiche dei due regimi:
\begin{itemize}
\item \textbf{Plasma classico:} Il parametro di accoppiamento scala come $\Gamma \propto n^{1/3} / T$. Questo significa che l'accoppiamento (e la non-idealità) \textbf{cresce} all'aumentare della densità.
\item \textbf{Plasma degenere:} Il parametro di accoppiamento quantistico scala come $r_s \propto n^{-1/3}$. In questo caso, l'accoppiamento \textbf{diminuisce} all'aumentare della densità, rendendo il gas via via più ideale sotto schiacciamento.
\item \textbf{Il caso delle nane bianche:} In un ambiente estremo come l'interno di una nana bianca, il plasma si sdoppia. Il fluido degli elettroni risulta fortemente degenere e si comporta come un gas di Fermi quasi ideale. Contemporaneamente, il fluido formato dagli ioni (non essendo degenere per via dell'enorme massa rispetto all'elettrone) soggiace alle regole della fisica classica e risulta essere \textbf{fortemente accoppiato}.
\end{itemize}

\section{Il Diagramma T-n (Temperatura-Densità)}

La comprensione complessiva dei vari stati della materia ionizzata si sintetizza elegantemente nel \textbf{diagramma T-n}, un grafico logaritmico che mappa la temperatura (asse delle ordinate) contro la densità elettronica (asse delle ascisse). All'interno di questo spazio delle fasi, possiamo tracciare diverse curve limitanti che suddividono i domini fisici, delineando i territori in cui si collocano i principali plasmi astrofisici e di laboratorio:

\begin{itemize}
\item \textbf{La soglia di Degenerazione (Linea Arancione, $T = T_F$):} Questa retta inclinata nel piano bi-logaritmico separa nettamente il plasma \textbf{classico} (che giace nella regione in alto a sinistra, a temperature elevate e basse densità) dal plasma \textbf{degenere} (localizzato in basso a destra, a basse temperature e altissime densità). Plasmi come il mezzo intergalattico (ICM), il vento solare e la corona risiedono saldamente nel dominio classico. L'interno di una nana bianca, al contrario, si colloca abbondantemente oltre questa linea, nella pura zona di degenerazione. Il nucleo solare (solar core) si posiziona esattamente in prossimità di questo confine, trovandosi in uno stato di "lieve degenerazione". L'energia termica incontra il limite intrinseco imposto dalla massa a riposo dell'elettrone ($k_B T = m_e c^2$), oltre il quale intervengono effetti relativistici estremi (segnato dalla linea tratteggiata orizzontale in alto).
\item \textbf{La transizione di Accoppiamento Forte (Linea Blu e curva, $\Gamma = 1$ e $r_s = 1$):} La linea curva (blu) traccia il luogo dei punti in cui l'energia termica classica eguaglia l'energia Coulombiana ($\Gamma = 1$), separando il plasma debolmente accoppiato ideale dal plasma fortemente accoppiato. Una volta intersecata la linea di degenerazione, questo limite si trasforma in una retta verticale ($r_s = 1$), poiché nel regime degenere l'accoppiamento cessa di dipendere dalla temperatura termica e dipende esclusivamente dalla densità. Le fotosfere stellari fredde, ad esempio, cadono nella regione di forte accoppiamento classico, mentre le nane bianche si pongono nella zona debolmente accoppiata degenere.
\item \textbf{Equilibrio di Ionizzazione (Curva Verde Acqua):} Questa curva rappresenta il luogo geometrico in cui il gas di idrogeno si presenta ionizzato al 50% all'equilibrio termodinamico locale (LTE), tracciato seguendo l'equazione di Saha. Le regioni al di sotto di tale linea non possiedono energia termica sufficiente per sostenere lo stato di plasma ideale.
\end{itemize}

\section{Analisi Comparata di Due Plasmi Densi: Nucleo Solare e Nana Bianca}

L'applicazione pratica dei concetti termodinamici appena esposti trova il suo massimo compimento nel confronto diretto tra le proprietà microscopiche di due plasmi stellari estremamente densi: il cuore termonucleare del Sole e l'interno di una stella degenere.

\subsection{Il Nucleo Solare (The Solar Core)}
All'interno del Sole, la densità elettronica stazionata a $n_e = 6 \times 10^{25} \text{ cm}^{-3}$ e la temperatura sfiora $T = 1.5 \times 10^7 \text{ K}$.
Sostituendo questi dati, il calcolo del parametro di accoppiamento classico restituisce $\Gamma \approx 0.07$, e il numero di particelle per sfera di Debye è circa $N_D \approx 10$. Di conseguenza, il plasma solare è \textbf{debolmente accoppiato, ma non in maniera trascurabile} ("weakly but not very weakly coupled").
Controllando la soglia quantistica, il rapporto $T/T_F$ è circa $2$, e il volume dell'onda termica $n\lambda_{\text{th}}^3 \approx 0.4$. Ciò indica che gli elettroni si trovano in una condizione di \textbf{lieve degenerazione} ("mildly degenerate"). Questa zona liminale produce delle discrepanze teoriche: le deviazioni dall'equazione di stato dei gas perfetti non sono astronomiche, ma si attestano su correzioni percentuali. Tuttavia, l'attuale strumentazione di \textbf{eliosismologia} è così raffinata da poter misurare e confermare empíricamente l'esistenza di queste correzioni termodinamiche di ordine percentuale.

\subsection{L'Interno di una Nana Bianca (A White-Dwarf Interior)}
In una tipica nana bianca a base carbonio-ossigeno, la densità esplode a $n_e \sim 10^{30} \text{ cm}^{-3}$, mantenendo una temperatura residua di $T \sim 10^7 \text{ K}$.
In questo ambiente, occorre scindere lo studio per specie:
\begin{itemize}
\item \textbf{Componente Elettronica:} Il rapporto $T/T_F$ crolla vertiginosamente a $\sim 10^{-3}$, mentre il parametro di accoppiamento quantistico $r_s \approx 0.01$. Gli elettroni costituiscono quindi un fluido \textbf{fortemente degenere e quasi perfettamente ideale}. La fortissima compressione ne annulla quasi del tutto la carica di accoppiamento, in accordo con la teoria.
\item \textbf{Componente Ionica (Ioni Carbonio):} Gli ioni massicci seguono invece la statistica classica. Calcolando l'accoppiamento ionico si trova $\Gamma_i = Z^2 e^2 / a_i T \approx 50$. L'alta densità comprime il nucleo classico spingendolo in regime di \textbf{forte accoppiamento} ("strongly coupled").
\end{itemize}
L'aspetto termodinamico più affascinante dell'evoluzione stellare insorge durante la senescenza termica della nana bianca. Via via che l'astro rilascia calore e si raffredda nello spazio vuoto, il denominatore $T$ si abbassa, facendo impennare inesorabilmente il parametro $\Gamma_i$. Quando il raffreddamento fa raggiungere la soglia critica di $\Gamma_i \to 175$, la teoria prevede che il nucleo ionico subisca una violenta transizione di fase e \textbf{cristallizzi}, formando un gigantesco reticolo metallico ordinato a livello macroscopico. Il rilascio repentino del calore latente di solidificazione si immette nel flusso termico e \textbf{ritarda visibilmente} il tasso di raffreddamento della stella, un effetto chiaramente osservabile nelle sequenze evolutive contemporanee.

\subsection{Evidenze Osservative della Cristallizzazione Stellare}
La predizione teorica, formulata negli anni '60 sulla base della termodinamica dei plasmi fortemente accoppiati, ha trovato negli ultimi anni una magnifica conferma empirica, resa possibile dall'altissima mole di dati fotometrici forniti dal telescopio spaziale astrometrico Gaia.

Tracciando un campione di 15.000 nane bianche collocate entro 100 parsec di distanza dal Sistema Solare, è stato possibile costruire uno straordinario diagramma colore-magnitudine, mettendo a confronto il colore astrometrico (Gaia BP-RP colour, misuratore indiretto della temperatura) con la magnitudine assoluta (Gaia G).
Guardando il grafico di dispersione della popolazione stellare in esame, invece di osservare uno sciame uniforme che declina uniformemente, appare in modo netto ed evidente un denso "accumulo" (\textit{pile-up}) di nane bianche raggruppate lungo una precisa fascia trasversale (banda evidenziata in arancione).

Questa striscia di densità anomala coincide inequivocabilmente con il percorso tracciato dal modello evolutivo per la sequenza di raffreddamento al momento della transizione. In questa determinata regione le stelle subiscono la cristallizzazione del nucleo di carbonio ($\Gamma_i \approx 175$). Il rilascio di \textbf{calore latente} compensa la perdita energetica, forzando la stella a stazionare (linger) alla medesima luminosità per un tempo molto prolungato, ritardando l'offuscamento ("cooling"). L'addensamento osservativo in questo diagramma colore-magnitudine costituisce dunque la riprova osservazionale diretta di una transizione di fase all'interno di un plasma fortemente accoppiato.

\section{L'Equilibrio di Ionizzazione: L'Equazione di Saha}

L'analisi dell'equilibrio di ionizzazione costituisce un pilastro fondamentale per comprendere lo stato termodinamico di un plasma astrofisico o di laboratorio. Il formalismo che descrive la ripartizione tra atomi neutri e ioni è racchiuso nell'Equazione di Saha.

\subsection{Un Paradosso Apparente: Il Ruolo dell'Entropia}
Prima di derivare rigorosamente le equazioni, è istruttivo affrontare un paradosso concettuale. Consideriamo l'atomo di idrogeno: la sua energia di ionizzazione è $I = 13.6 \text{ eV}$. Convertendo questo valore in una temperatura equivalente, dividendo per la costante di Boltzmann ($k_B$), si ottiene $I/k_B = 158\,000 \text{ K}$.

Eppure, l'osservazione empirica mostra uno scenario apparentemente in contrasto con questo altissimo valore teorico:
\begin{itemize}
\item La fotosfera solare, che si trova a "sole" $5800 \text{ K}$, è debolmente ionizzata.
\item Nelle atmosfere stellari, l'idrogeno risulta già \textbf{mezzo ionizzato} a temperature modeste di $\approx 10^4 \text{ K}$.
\item Nell'Universo primordiale (epoca della ricombinazione), la transizione di ionizzazione è avvenuta a temperature ancora più basse, attorno ai $\approx 3700 \text{ K}$.
\end{itemize}

Come è possibile che il gas si ionizzi a temperature di oltre un ordine di grandezza inferiori rispetto all'energia di legame $I$? Il \textbf{risultato chiave} per risolvere questo enigma risiede nella termodinamica statistica: \textbf{l'entropia vince}.
Mentre un elettrone legato al nucleo occupa un singolo stato quantico ben definito, un elettrone libero può occupare una vastissima gamma di stati accessibili nello spazio delle fasi, proporzionale a $\sim 1/(n\lambda_{\text{th}}^3) \gg 1$, dove $\lambda_{\text{th}}$ è la lunghezza d'onda termica di de Broglie e $n$ la densità.

La transizione verso lo stato ionizzato diventa statisticamente e termodinamicamente favorita non appena il peso del fattore di Boltzmann eguaglia il numero di stati liberi disponibili, ovvero quando:
\begin{equation}
e^{-I/T} \gtrsim n\lambda_{\text{th}}^3 \implies T_{\text{ion}} \approx \frac{I}{\ln(1/n\lambda_{\text{th}}^3)} \sim \frac{I}{20 \div 40}
\end{equation}
La conseguenza fisica diretta di questa relazione è che \textbf{più il gas è diluito (bassa densità $n$), minore sarà la temperatura di ionizzazione necessaria} per sottrarre l'elettrone al nucleo.

\subsection{Derivazione (1): Gas Ideale con Stati Interni}
Per formalizzare matematicamente questo equilibrio, partiamo dalla funzione di partizione per un gas ideale le cui particelle possiedono gradi di libertà interni. Consideriamo una singola particella appartenente alla specie $s$. La sua funzione di partizione di singola particella, $Z_{1,s}$, è il prodotto del contributo traslazionale e degli stati interni (caratterizzati dal peso statistico $U_s$ e dall'energia dello stato fondamentale $E_s$):
\begin{equation}
Z_{1,s} = \underbrace{\int \frac{d^3x d^3p}{h^3} e^{-p^2/2m_sT}}_{V (2\pi m_s T)^{3/2} / h^3} \times U_s e^{-E_s/T} = \frac{V U_s}{\lambda_s^3} e^{-E_s/T}
\end{equation}
dove abbiamo introdotto la lunghezza d'onda termica $\lambda_s = \frac{h}{(2\pi m_s T)^{1/2}}$.

Passando a un sistema macroscopico costituito da $N_s$ particelle identiche e non interagenti, la funzione di partizione totale $Z$ si ottiene dividendo per $N_s!$ per correggere l'indistinguibilità: $Z = Z_{1,s}^{N_s}/N_s!$.
Per grandi numeri, applichiamo l'approssimazione di Stirling, $\ln N! \simeq N\ln N - N$, e ricaviamo l'energia libera di Helmholtz $F$:
\begin{equation}
F = -T \ln Z = -N_s T \left[ \ln \left( \frac{Z_{1,s}}{N_s} \right) + 1 \right]
\end{equation}

Dall'energia libera possiamo ricavare il \textbf{potenziale chimico} $\mu_s$, definito come la derivata parziale di $F$ rispetto al numero di particelle a temperatura e volume fissati:
\begin{equation}
\mu_s = \left. \frac{\partial F}{\partial N_s} \right|*{T,V} = -T \ln \left( \frac{Z*{1,s}}{N_s} \right)
\end{equation}
Invertendo quest'ultima equazione e sostituendo l'espressione di $Z_{1,s}$, isoliamo la densità numerica di particelle della specie $s$ ($n_s = N_s/V$):
\begin{equation}
n_s = \frac{N_s}{V} = U_s \left( \frac{2\pi m_s T}{h^2} \right)^{3/2} e^{(\mu_s - E_s)/T}
\end{equation}

\subsection{Derivazione (2): Equilibrio Chimico}
Consideriamo ora la reazione reversibile di ionizzazione dell'idrogeno: $\text{H} \rightleftharpoons p + e$ (un atomo di idrogeno si dissocia in un protone e un elettrone).
All'equilibrio termodinamico, mantenendo $T$ e $V$ costanti, l'energia libera $F$ deve raggiungere un \textbf{minimo} ($dF = 0$). Data la stechiometria della reazione, le variazioni del numero di particelle sono vincolate dalla relazione $dN_p = dN_e = -dN_{\text{H}}$. Sostituendo nel differenziale totale dell'energia libera otteniamo:
\begin{equation}
dF = \mu_p dN_p + \mu_e dN_e + \mu_{\text{H}} dN_{\text{H}} = 0 \implies \mu_{\text{H}} = \mu_p + \mu_e
\end{equation}
Questa è la condizione di equilibrio chimico. Formiamo adesso il rapporto delle densità $n_p n_e / n_{\text{H}}$, utilizzando per ciascuna specie la formula per la densità $n_s$ ricavata nel passaggio precedente:
\begin{equation}
\frac{n_p n_e}{n_{\text{H}}} = \frac{U_p U_e}{U_{\text{H}}} \left( \frac{2\pi T}{h^2} \right)^{3/2} \left( \frac{m_p m_e}{m_{\text{H}}} \right)^{3/2} e^{(\mu_p + \mu_e - \mu_{\text{H}})/T} e^{-(E_p + E_e - E_{\text{H}})/T}
\end{equation}

Questa complessa espressione si semplifica drasticamente grazie a tre considerazioni fisiche fondamentali:
\begin{itemize}
\item La somma algebrica dei potenziali chimici all'esponente è identicamente nulla in virtù della condizione di equilibrio ($\mu_p + \mu_e - \mu_{\text{H}} = 0$). I potenziali chimici incogniti si elidono del tutto.
\item La massa del protone è pressocché identica alla massa dell'atomo di idrogeno intero ($m_p \simeq m_{\text{H}}$). Il loro rapporto si approssima a 1, e nel termine di massa sopravvive unicamente il contributo elettronico ($m_e^{3/2}$).
\item La differenza di energia fondamentale tra i prodotti isolati e l'atomo legato corrisponde esattamente all'energia di ionizzazione: $E_p + E_e - E_{\text{H}} = I$.
\end{itemize}

\subsection{L'Equazione di Saha}
Applicando le approssimazioni discusse, il rapporto all'equilibrio per l'idrogeno si condensa nell'espressione:
\begin{equation}
\frac{n_p n_e}{n_{\text{H}}} = \frac{U_p U_e}{U_{\text{H}}} \left( \frac{2\pi m_e T}{h^2} \right)^{3/2} e^{-I/T}
\end{equation}
Per calcolare esattamente il pre-fattore, analizziamo i pesi statistici:
\begin{itemize}
\item Per l'elettrone libero, $U_e = 2$ (dovuto ai due possibili stati di spin).
\item Per il protone nudo, convenzionalmente $U_p = 1$.
\item Per l'idrogeno neutro nello stato fondamentale (1s), si ha $U_{\text{H}} \simeq 2$ (vi sono due stati di spin dell'elettrone associato; il contributo dello spin nucleare si cancella nel bilancio complessivo).
\end{itemize}

Generalizzando questo \textbf{risultato chiave} (formulato da Meghnad Saha nel 1920) a \textit{qualsiasi} elemento chimico che passi da uno stadio di ionizzazione $i$ allo stadio successivo $i+1$, otteniamo la celebre \textbf{Equazione di Saha}:
\begin{equation}
\frac{n_{i+1} n_e}{n_i} = 2 \frac{U_{i+1}}{U_i} \left( \frac{2\pi m_e T}{h^2} \right)^{3/2} e^{-I_i/T} \equiv K_i(T)
\end{equation}
Punti focali di questo risultato:
\begin{itemize}
\item Il fattore $(2\pi m_e T/h^2)^{3/2}$ corrisponde a $\lambda_{\text{th}}^{-3}$, ovvero la \textbf{concentrazione quantistica degli elettroni liberi}. La sua comparsa nell'equazione rende quantitativamente esatta la risoluzione del paradosso entropico introdotto inizialmente.
\item Si deve scrivere e risolvere una diversa Equazione di Saha per \textit{ciascuno} stadio di ionizzazione (es. da neutro a singolarmente ionizzato, da singolarmente a doppiamente ionizzato, ecc.).
\item La validità di questo modello postula rigorosamente la sussistenza dell'\textbf{Equilibrio Termodinamico Locale (LTE)} e presuppone che le particelle si comportino come un \textbf{gas ideale}.
\end{itemize}

\subsection{Il Caso Pratico dell'Idrogeno Puro}
Mettiamo in pratica l'equazione studiando un gas costituito interamente da idrogeno puro.
Definiamo la densità totale di nuclei (sia legati che liberi) come $n = n_{\text{H}} + n_p$. Introduciamo la \textbf{frazione di ionizzazione} $x$, definita come $x = n_p / n$. Affinché il plasma mantenga la neutralità macroscopica di carica, la densità di elettroni deve pareggiare quella dei protoni liberi: $n_e = n_p$.

Con queste variabili, possiamo riscrivere le singole densità in funzione di $n$ e $x$:
\begin{equation}
n_p = n_e = xn, \quad \quad n_{\text{H}} = (1 - x)n
\end{equation}
Sostituendo tali espressioni nell'equazione di Saha si ha:
\begin{equation}
\frac{(xn)^2}{(1 - x)n} = K(T) \implies \frac{x^2}{1 - x} = \frac{K(T)}{n} = \frac{1}{n \lambda_{\text{th}}^3} e^{-I/T}
\end{equation}

Inserendo i valori numerici delle costanti fisiche fondamentali, la concentrazione quantistica fornisce il prefattore $(2\pi m_e k_B / h^2)^{3/2} = 2.415 \times 10^{15} \text{ cm}^{-3} \text{ K}^{-3/2}$, mentre l'energia di ionizzazione normalizzata vale $I/k_B = 157\,900 \text{ K}$.

Sostituendo queste grandezze, approdiamo al \textbf{risultato chiave} operativo per calcolare la frazione di ionizzazione dell'idrogeno:
\begin{equation}
\frac{x^2}{1 - x} = \frac{2.415 \times 10^{15}}{n , [\text{cm}^{-3}]} T_{\text{K}}^{3/2} \exp\left( -\frac{157,900}{T_{\text{K}}} \right)
\end{equation}
dove $T_{\text{K}}$ è la temperatura espressa in Kelvin e $n$ la densità in particelle al centimetro cubo.

\section{Risoluzione dell'Equazione di Saha}

L'equazione di Saha, che mette in relazione la frazione di ionizzazione $x$ con la temperatura $T$ e la densità totale $n$, può essere riscritta definendo il parametro aggregato $r = K/n$. L'equazione assume così la forma algebrica:
\begin{equation}
x^2 + rx - r = 0
\end{equation}
. Risolvendo questa equazione quadratica per $x$, e scartando la radice negativa in quanto priva di significato fisico, la soluzione per la frazione di ionizzazione risulta:
\begin{equation}
x = \frac{-r + \sqrt{r^2 + 4r}}{2} = \frac{2}{1 + \sqrt{1 + 4/r}}
\end{equation}
. La seconda forma proposta per questa soluzione è numericamente vantaggiosa poiché evita la cancellazione numerica quando $r \gg 1$ (ovvero nel limite di forte ionizzazione).

\subsection{La Temperatura di Mezza Ionizzazione ($T_{1/2}$)}
Un parametro fisico di grande interesse per caratterizzare il plasma è la temperatura alla quale esattamente la metà degli atomi risulta ionizzata, definita come \textbf{temperatura di mezza ionizzazione} ($T_{1/2}$). In questa condizione, la frazione di ionizzazione è per definizione $x = 1/2$. Sostituendo questo valore nell'equazione quadratica, si deduce immediatamente che $r = 1/2$, il che implica a sua volta $n = 2K(T_{1/2})$.
Applicando il logaritmo naturale ad ambo i membri e isolando $T_{1/2}$, si ottiene:
\begin{equation}
\ln n = \ln \left( 2 \times 2.415 \times 10^{15} T_{1/2}^{3/2} \right) - \frac{I}{T_{1/2}} \implies T_{1/2} = \frac{I}{\ln \left[ 4.83 \times 10^{15} T_{1/2}^{3/2} / n \right]}
\end{equation}
.
Questa equazione trascendente per $T_{1/2}$ rivela diverse peculiarità:
\begin{itemize}
\item Può essere facilmente risolta per iterazione numerica. Partendo da un valore di primo tentativo, ad esempio $T_{1/2} = 10^4 \text{ K}$, la convergenza verso il valore esatto avviene in appena $2-3$ passaggi.
\item Il rapporto $I/T_{1/2}$ coincide con il denominatore logaritmico, il cui argomento è un numero enorme, per cui il valore tipico si aggira nell'intervallo $\approx 15-50$.
\item L'aspetto termodinamico più rilevante è che **minore è la densità del gas $n$, maggiore sarà l'argomento del logaritmo**, con il risultato che la temperatura di mezza ionizzazione $T_{1/2}$ si abbassa progressivamente nei mezzi rarefatti.
\end{itemize}

\subsection{La Ripidezza della Transizione (Sharpness)}
È lecito domandarsi con quale rapidità termica si compia il passaggio dallo stato neutro a quello ionizzato. Per quantificare l'estensione dell'intervallo di transizione, si differenzia la relazione logaritmica dell'equazione di Saha, $\ln[x^2/(1-x)] = \ln r$, tenendo conto della dipendenza $r \propto T^{3/2}e^{-I/T}$:
\begin{equation}
\left( \frac{2}{x} + \frac{1}{1-x} \right) dx = d\ln r = \left( \frac{3}{2} + \frac{I}{T} \right) d\ln T
\end{equation}
.
Valutiamo questa espressione nel punto di mezza ionizzazione ($x = 1/2$). Il termine in parentesi a sinistra assume il valore $4 + 2 = 6$. Isolando la derivata della frazione di ionizzazione rispetto al logaritmo della temperatura, si ha:
\begin{equation}
\left. \frac{dx}{d\ln T} \right|*{x=1/2} = \frac{1}{6} \left( \frac{3}{2} + \frac{I}{T*{1/2}} \right) \approx 3-8
\end{equation}
.
Questo \textbf{risultato chiave} dimostra che il processo non è affatto graduale: la frazione di ionizzazione passa dal 10% ($x=0.1$) al 90% ($x=0.9$) su un intervallo di temperatura ristrettissimo, pari a una variazione relativa $\Delta T / T \approx 10\%-20\%$. La ionizzazione è pertanto \textbf{netta e improvvisa} (\textit{sharp}) perché la temperatura interviene all'esponente di un termine dominato da un coefficiente molto grande, $I/T$.

Questa "violenza" della transizione è chiaramente visibile tracciando le **curve di ionizzazione dell'idrogeno** (frazione $x$ in funzione di $T$ per diversi valori costanti di densità totale $n_{\text{H}} + n_p$). Al diminuire della densità (da $10^{20}$ fino a $1 \text{ cm}^{-3}$), le curve si spostano rigidamente verso sinistra, mantenendo inalterato il loro profilo a gradino molto ripido. Tutte le transizioni termiche calcolate avvengono a temperature confinate molto al di sotto del valore di ionizzazione asintotico $I/k_B = 1.58 \times 10^5 \text{ K}$. A titolo esemplificativo, alla densità della fotosfera $10^{15} \text{ cm}^{-3}$, la $T_{1/2}$ si aggira a $10\,200 \text{ K}$, mentre in condizioni di estrema diluizione ($1 \text{ cm}^{-3}$), essa crolla a circa $3300 \text{ K}$.

\section{L'Equazione di Saha in Ambienti Astrofisici Reali}

Per concludere l'analisi dell'equilibrio di ionizzazione, si esaminano due scenari astrofisici in cui l'equazione di Saha svolge un ruolo critico nel determinare le proprietà radiative del plasma.

\subsection{La Fotosfera Solare}
Nella regione atmosferica corrispondente alla fotosfera solare, i parametri fisici si assestano su una densità $n_{\text{H}} \approx 10^{17} \text{ cm}^{-3}$ e una temperatura di superficie $T = 5800 \text{ K}$. Applicando rigorosamente l'equazione di Saha per l'idrogeno, si ricava una frazione di ionizzazione estremamente bassa, $x_{\text{H}} \approx 10^{-4}$.

Tuttavia, l'apparente contraddizione di un plasma con così pochi elettroni liberi viene risolta introducendo la fisica degli elementi "contaminanti". Nella fotosfera superiore abbondano atomi metallici pesanti (quali Na, Mg, Fe) caratterizzati da potenziali di prima ionizzazione molto deboli, compresi nell'intervallo $I \approx 5-8 \text{ eV}$. L'applicazione di Saha a questi elementi dimostra che sono loro a \textbf{fornire la quasi totalità degli elettroni liberi} nel mezzo.
L'eccesso di elettroni liberi metallici induce la formazione dell'esotico ione idrogeno negativo ($\text{H}^-$), il quale funge da \textbf{fonte primaria di opacità continua} (continuum opacity) nella banda ottica del visibile, impedendoci di osservare gli strati solari più profondi.

\subsection{Ricombinazione Cosmologica (Cosmological Recombination)}
Traslando il modello ai primordi dell'Universo, l'equazione di Saha governa la Ricombinazione (o disaccoppiamento della materia dalla radiazione).
Nel gas primordiale, il rapporto densità barionica su fotoni si attesta a un valore irrisorio, $\eta = n_b / n_\gamma \approx 6 \times 10^{-10}$, implicando una "minuscola densità effettiva" (a tiny effective density) per l'interazione elettrone-protone.

In virtù di questa rarefazione estrema (come dimostrato in precedenza dalla forte dipendenza dalla densità del logaritmo), Saha prevede che la mezza ionizzazione ($x=1/2$) venga raggiunta a temperature ridotte, vicine a $T \approx 3700 \text{ K}$, a un redshift cosmologico $z \approx 1400$. Nonostante la discesa al di sotto di questa temperatura teorica di Saha proceda spedita con l'espansione, complessi effetti radiativi sub-atomici (principalmente la fotoionizzazione causata dall'emissione intrappolata della linea Lyman-$\alpha$) ritardano l'ultimo scattering fotone-materia fino a redshift inferiori, cristallizzando la Radiazione Cosmica di Fondo a $z \approx 1100$ (come dimostrato dalle stime pionieristiche di Peebles, 1968, e successivamente confermato dalle mappe del satellite Planck dell'ESA).

\section{I Limiti dell'Equazione di Saha (Parte I): La Divergenza della Funzione di Partizione}

Un primo limite intrinseco nell'applicazione rigorosa dell'equazione di Saha riguarda il calcolo della funzione di partizione, la quale tende matematicamente a divergere all'infinito. Considerando il generico livello energetico $n$ dell'atomo di idrogeno, la sua energia di legame è data da $-I/n^2$, mentre il suo peso statistico ammonta a $2n^2$. La funzione di partizione per l'idrogeno neutro ($U_{\text{H}}$) si esprime quindi come la somma su tutti i possibili stati legati:
\begin{equation}
U_{\text{H}} = \sum_{n=1}^\infty 2n^2 e^{I/n^2 T}
\end{equation}
Calcolando il limite per $n \to \infty$, il termine esponenziale tende asintoticamente a $1$, e la serie si riduce alla somma dei termini $2n^2$, che diverge inevitabilmente a infinito ($\sum 2n^2 = \infty$).

Fisicamente, questo paradosso matematico scaturisce dall'assunzione errata che un atomo isolato possieda un numero infinito di stati legati, i cui raggi orbitali crescono proporzionalmente a $n^2 a_B$ (dove $a_B$ è il raggio di Bohr). All'interno di un plasma reale, tuttavia, l'atomo non è mai perfettamente isolato: le orbite elettroniche la cui estensione supera il valore minimo tra la distanza media interparticellare $a$ e la lunghezza di schermo di Debye $\lambda_D$ vengono letteralmente distrutte dalle perturbazioni introdotte dalle cariche vicine. Di conseguenza, è necessario troncare fisicamente la sommatoria a un livello massimo $n_{\text{max}}$, definito in modo tale che $n_{\text{max}}^2 a_B \sim \min(a, \lambda_D)$. A basse densità, questo troncamento risulta marginale e l'approssimazione dello stato fondamentale $U_{\text{H}} \simeq 2$ si rivela adeguata; ad alte densità, invece, il troncamento altera profondamente il bilancio termodinamico.

\section{L'Abbassamento del Continuo (Continuum Lowering)}

Un'ulteriore e severa limitazione dell'equazione di Saha nei plasmi densi è rappresentata dal fenomeno noto come abbassamento del continuo. In prossimità di uno ione avente carica $Ze$, il potenziale elettrico non assume la classica forma coulombiana, ma risulta schermato dal plasma circostante. Per distanze molto inferiori alla lunghezza di Debye ($r \ll \lambda_D$), il potenziale schermato può essere approssimato espandendo in serie di Taylor la componente esponenziale:
\begin{equation}
\phi(r) = \frac{Ze}{r} e^{-r/\lambda_D} \simeq \frac{Ze}{r} \left( 1 - \frac{r}{\lambda_D} \right) = \frac{Ze}{r} - \frac{Ze}{\lambda_D}
\end{equation}
Questa espressione analitica rivela che l'elettrone legato si trova all'interno di una buca di potenziale Coulombiana che risulta uniformemente "sollevata" (shifted up) di una quantità pari a $Ze/\lambda_D$. A causa di questo innalzamento strutturale del livello di base, l'elettrone raggiunge lo stato libero del continuo a un livello energetico di $-Ze^2/\lambda_D$ anziché al consueto valore $0$.

Il \textbf{risultato chiave} che discende da questa dinamica è che l'energia di ionizzazione effettiva del sistema ($I_{\text{eff}}$) risulta essere sensibilmente inferiore rispetto a quella dell'atomo isolato ($I$):
\begin{equation}
I_{\text{eff}} = I - \Delta I, \quad \Delta I \simeq \frac{Ze^2}{\lambda_D}
\end{equation}
Questa formulazione matematica corrisponde esplicitamente al limite di bassa densità descritto dalla teoria di Stewart-Pyatt. Valutando il rapporto adimensionale $\Delta I / T = Z e^2 / \lambda_D T \sim 1/N_D$, si evince in modo chiaro che tale correzione energetica è del tutto trascurabile all'interno dei plasmi ideali (poiché $N_D \gg 1$), ma diventa imponente e impossibile da ignorare nei plasmi ad altissima densità.

\section{Il Fallimento di Saha nel Nucleo Solare}

L'applicazione ingenua (naive) dell'equazione di Saha standard alle condizioni del nucleo solare fornisce un esempio emblematico dei limiti teorici appena esposti. Inserendo all'interno dell'equazione i parametri estremi del nucleo, ovvero una densità elettronica $n = 6 \times 10^{25} \text{ cm}^{-3}$ e una temperatura di $T = 1.5 \times 10^7 \text{ K}$, il calcolo prevede una frazione di ionizzazione $x \approx 0.75$. Questo risultato paradossale implicherebbe che circa un quarto dell'idrogeno stellare si trovi ancora allo stato neutro all'interno del nucleo termonucleare, un dato palesemente inesatto.

Le ragioni fisiche di questo vistoso fallimento risiedono in tre fattori ambientali dominanti:
\begin{itemize}
\item La distanza media interparticellare all'interno del plasma solare, calcolata in $a \approx 1.6 \times 10^{-9} \text{ cm}$, risulta essere strettamente inferiore al raggio di Bohr atomico $a_B = 5.3 \times 10^{-9} \text{ cm}$, proibendo fisicamente la stessa sussistenza geometrica dell'orbita legata dell'idrogeno.
\item L'impiego rigoroso del potenziale di schermatura di Yukawa ($e^{-r/\lambda_D}/r$) dimostra matematicamente l'impossibilità di supportare stati legati stabili qualora la lunghezza di Debye sia ridotta sotto la soglia critica di $\lambda_D < 0.84 a_B$. Nel nucleo solare, la sola lunghezza di Debye elettronica si ferma a $\lambda_{De} \approx 0.65 a_B$, e precipita a $0.41 a_B$ quando si include la reazione di schermo degli ioni.
\item L'effetto di abbassamento del continuo ($\Delta I = e^2/\lambda_D$) calcolato per queste densità ammonta a circa $65 \text{ eV}$, una correzione enormemente superiore alla stessa energia di ionizzazione di base dell'idrogeno ($I = 13.6 \text{ eV}$).
\end{itemize}
In virtù di queste evidenze, la conclusione teorica corretta è che il nucleo stellare risulta essere totalmente ionizzato a causa della \textbf{ionizzazione per pressione} (pressure ionisation), un meccanismo quantistico governato dalla compressione spaziale che la classica equazione di Saha, basata sul limite di gas ideale, è strutturalmente incapace di prevedere o descrivere.

\section{Limiti dell'Equazione di Saha (Parte II): Deviazioni dall'Equilibrio Termodinamico Locale (LTE)}

In tutti i regimi in cui il plasma non soddisfa le rigide assunzioni dell'Equilibrio Termodinamico Locale (LTE), l'uso dell'equazione di Saha perde validità, imponendo un approccio cinetico differenziale in cui l'evoluzione temporale della densità elettronica viene tracciata computando individualmente ogni singolo processo elementare. L'equazione di bilancio assume quindi la forma:
\begin{equation}
\frac{dn_e}{dt} = \underbrace{n_e n_{\text{H}} S(T)}*{\text{ionizz. collisionale}} - \underbrace{n_e^2 n_p \beta(T)}*{\text{ricomb. a 3 corpi}} + \underbrace{n_{\text{H}} \Gamma_{\text{ph}}}*{\text{foto-ionizz.}} - \underbrace{n_e n_p \alpha*{\text{rr}}(T)}_{\text{ricomb. radiativa}}
\end{equation}
.

Solamente sotto le condizioni di LTE si attua il principio del bilancio dettagliato (detailed balance): ogni specifico processo elementare deve bilanciare quantitativamente in maniera esatta e indipendente il proprio processo inverso. Uguagliando, ad esempio, i tassi assoluti di ionizzazione collisionale e ricombinazione a tre corpi (3-body recombination), si ottiene un'elegante via per derivare coefficienti complessi:
\begin{equation}
n_e n_{\text{H}} S = n_e^2 n_p \beta \implies \frac{n_p n_e}{n_{\text{H}}} = \frac{S}{\beta} = K(T) \implies \beta(T) = \frac{S(T)}{K(T)}
\end{equation}
Questa concatenazione logica dimostra come il calcolo di un rateo di interazione "difficile" (come la ricombinazione a tre corpi, definita $\beta(T)$) possa essere ricavato in maniera indiretta conoscendo il corrispondente tasso elementare di collisione $S(T)$ e combinandolo opportunamente con la costante di equilibrio di Saha $K(T)$. Allo stesso identico modo, il rateo di foto-ionizzazione $\Gamma_{\text{ph}}$, immerso all'interno di un intenso campo radiativo di corpo nero, si porrà in perfetto e rigoroso equilibrio rispetto al tasso di ricombinazione radiativa $\alpha_{\text{rr}}$.

\section{L'Equilibrio Coronale (Coronal Equilibrium)}

Quando l'ambiente plasmatico si presenta contestualmente iper-riscaldato, estremamente diluito e otticamente sottile (optically thin) — condizioni tipiche della corona solare — l'equilibrio dei tassi subisce una radicale alterazione. La drastica bassa densità e l'incapacità di confinare la radiazione causano la disattivazione sia della foto-ionizzazione sia della complessa ricombinazione a tre corpi. Di conseguenza, lo stato di ionizzazione del sistema è mantenuto da un esclusivo bilancio tra la ionizzazione collisionale innescata dall'agitazione termica e la ricombinazione spontanea radiativa:
\begin{equation}
n_e n_{\text{H}} S = n_e n_p \alpha_{\text{rr}} \implies \frac{n_p}{n_{\text{H}}} = \frac{x}{1-x} = \frac{S(T)}{\alpha_{\text{rr}}(T)}
\end{equation}
Il \textbf{risultato chiave} di questa particolare equivalenza cinetica risiede nel fatto che la densità elettronica $n_e$ \textbf{si elide matematicamente (cancels out)} da ambedue i membri dell'equazione. Ne consegue che l'equilibrio di ionizzazione in ambiente coronale è univocamente dipendente dalla temperatura cinetica $T$, sganciandosi completamente dalla densità locale del plasma.

Questa esclusiva correlazione termica offre un ineguagliabile strumento di diagnostica astronomica, permettendo la calibrazione di precisissimi \textbf{termometri coronali}. Poiché ogni specifica configurazione ionica "sopravvive" stabilmente unicamente all'interno di un ristrettissimo intervallo (narrow range) di temperatura, l'isolamento visivo di specifiche righe spettrali ioniche costituisce una tomografia termica del Sole. Ad esempio, mappare il plasma inquadrando l'emissione del Fe IX a 171 Å rivela con esclusività le strutture a 0.7 MK, mentre inquadrando simultaneamente il Fe XII a 193 Å si mettono in risalto solo e unicamente i domini plasmatici riscaldati a 1.5 MK. Questo permette di osservare complessi strutturali del tutto indipendenti e differenti pur analizzando la stella nel medesimo preciso istante.

\section{Riepilogo delle Proprietà Fondamentali dello Stato di Plasma}

Volgendo al termine del modulo, procediamo con la sintesi delle metriche primarie dedotte che tracciano l'identikit formale dello stato di plasma.
\begin{itemize}
\item La \textbf{lunghezza di Debye} costituisce l'orizzonte di schermatura spaziale del gas: $\lambda_{De} = (T_e / 4\pi n_e e^2)^{1/2}$.
\item La \textbf{frequenza di plasma} scandisce il ritmo temporale delle oscillazioni collettive: $\omega_{pe} = (4\pi n_e e^2 / m_e)^{1/2}$.
\item Il loro rigoroso \textbf{prodotto} fissa il ponte cinematico alla velocità termica del mezzo: $\omega_{pe} \lambda_{De} = v_{te}$.
\item Il \textbf{parametro di plasma}, perno del regime di accoppiamento: $N_D = \frac{4\pi}{3} n \lambda_D^3 = 1 / (3\sqrt{3}\Gamma^{3/2})$.
\item Il termine descrittivo della \textbf{degenerazione quantistica} (Temperatura di Fermi): $T_F = \hbar^2 (3\pi^2 n_e)^{2/3} / 2m_e$.
\item La costante dell'\textbf{equazione di Saha} per i bilanci chimico-ionizzanti in gas ideale: $\frac{n_{i+1} n_e}{n_i} = 2 \frac{U_{i+1}}{U_i} \lambda_{\text{th}}^{-3} e^{-I_i/T}$.
\end{itemize}

\subsection{Criteri Generali e Sviluppi Futuri}
A titolo definitivo, affinché una porzione di materia ionizzata assuma le peculiarità di un plasma reale, essa deve inderogabilmente soddisfare la triade di disuguaglianze strutturali: il contenitore macroscopico deve sovrastare ogni fluttuazione di carica ($L \gg \lambda_D$), il campo schermante deve attingere a una ricca distribuzione di corpuscoli interattivi ($N_D \gg 1$), e la comunicazione elettrodinamica tra di essi deve sopravvivere senza venire dissipata dal trascinamento collisionale della materia neutra ($\omega_{pe} \tau_n \gg 1$).
A causa dell'esponenziale peso entropico dello spazio delle fasi quantistico, si osserva sistematicamente che la cascata di ionizzazione prende vita, con forte contrasto predittivo, già a temperature di gran lunga inferiori al potenziale isolato, aggirandosi nella sfera di $T \sim I/30$.

Il prosieguo accademico, demandato al Modulo 2a, sposterà l'indagine verso la meccanica fine delle \textbf{collisioni Coulombiane} a corto raggio, formulando l'incidenza del logaritmo di Coulomb ($\ln \Lambda$) nella resistività macroscopica e approfondendo le emissioni per \textit{bremsstrahlung} (radiazione di frenamento).